\section{USAJMO 2020}
        \begin{enumerate}
        	\item (\textit{USAJMO 2020}) \textbf{Note:} For any geometry problem whose statement begins with an asterisk $(*)$, the first page of the solution must be a large, in-scale, clearly labeled diagram. Failure to meet this requirement will result in an automatic 1-point deduction.


 Let $n \geq 2$ be an integer. Carl has $n$ books arranged on a bookshelf. Each book has a height and a width. No two books have the same height, and no two books have the same width. Initially, the books are arranged in increasing order of height from left to right. In a move, Carl picks any two adjacent books where the left book is wider and shorter than the right book, and swaps their locations. Carl does this repeatedly until no further moves are possible.  Prove that regardless of how Carl makes his moves, he must stop after a finite number of moves, and when he does stop, the books are sorted in increasing order of width from left to right.


 



\textbf{Solution}

Without loss of generality, let the books have widths and heights of 1, 2, 3,..., and n. First we will show that Carl can only make finitely many moves.

 Note that two books can never be swapped more than once with each other. This is because after a first swap, the wider book will be to the left of the other, so they can never be swapped back. Thus there can be at most ${n\choose 2}$ moves.

 Now we will show that the books are in increasing order of width when Carl is done making moves. Assume otherwise for the sake of contradiction. Now let $f(k)$ be the width of the book in the kth spot from the left. Since the books are not in increasing order of width, there must be i such that $f(i)>f(i+1)$, for $1\leq i\leq n-1$. 

 \textbf{Claim:} The ith book must be shorter than the $(i+1)$th book. 

 \textbf{Proof:} Assume otherwise for contradiction. Then the ith book must have started to the right of the $(i+1)$th book, since the ith book is taller. However, the ith book is now to the left of the $(i+1)$th book, so they must have been swapped at some point. But since $f(i)>f(i+1)$, the ith book is also wider. But the ith book started to the right, so since it was wider it cannot have been swapped with the $(i+1)$th book, so we arrive at a contradiction. $\blacksquare$

 From the claim, the ith book must be shorter than the $(i+1)$th book. But since the ith book is also wider by definition (as $f(i)>f(i+1)$), Carl can swap those two books, so he still has legal moves and we arrive at a contradiction, so we are done. $\blacksquare$

 


 


	\item (\textit{USAJMO 2020}) \textbf{Note:} For any geometry problem whose statement begins with an asterisk $(*)$, the first page of the solution must be a large, in-scale, clearly labeled diagram. Failure to meet this requirement will result in an automatic 1-point deduction.


 Let $\omega$ be the incircle of a fixed equilateral triangle $ABC$. Let $\ell$ be a variable line that is tangent to $\omega$ and meets the interior of segments $BC$ and $CA$ at points $P$ and $Q$, respectively. A point $R$ is chosen such that $PR = PA$ and $QR = QB$. Find all possible locations of the point $R$, over all choices of $\ell$.


 



\textbf{Solution 1}

Call a point \textit{good} if it is a possible location for $R$.Let the incircle of $\triangle ABC$ touch $BC$ at $D$, $AC$ at $E$, and $\ell$ at $T$. Also, let the center of the incircle be $I$. Clearly, a point is good iff it lies on the circle containing $A$ with center $P$ as well as the circle containing $B$ with center $Q$. Call these circles $\omega_1$ and $\omega_2$, respectively.

 Note that point $T$ can only lie on minor arc $\overarc{DE}$ (excluding the endpoints). 

 \textbf{Claim:} A point $X$ is good iff $XT\perp PQ$ and $XT=AD=BE$.

 \textbf{Proof:} WLOG, let $X$ be on the same side of $PQ$ as $I$. Then we have that $T, I, X$ are collinear. In particular, we have $\angle XTP=90^{\circ}$. Then $TX=AD$, $PT=PD$, $\angle XTP=\angle ADP$, so that $\triangle ADP\cong \triangle XTP\rightarrow PA=PX$. Similarly, $QB=QX$, so $X$ is good. Then $X'$, the reflection of $X$ across $T$, is also good. But $\omega_1$ and $\omega_2$ have at most two intersections, so $X$ and $X'$ must be these intersections, and since a point is good iff it lies on both circles, we are done.

 Now, we know that $R, I, T$ are collinear. Then we have two cases:

 \textbf{Case 1:} $R$ and $I$ lie on the same side of $PQ$. Then we have $RI+IT=AI+ID=BI+IE$, so that $RI=AI=BI$. Then we have that $R, A, B$ lie on a circle with center $I$. Note that because $T$ lies on $\overarc{DE}$, $R$ must lie on $\overarc{AB}$. So, one of the solutions is $\boxed{R \text{ in arc } AB\text{ (excluding the endpoints)}}$.

 \textbf{Case 2:} $R$ and $I$ lie on opposite sides of $PQ$. Then extend $ID$ out from $D$ to $X$ such that $DX=TR$, and extend $IE$ out from $E$ to $Y$ such that $EY=TR$. Then we see that $IR=IX=IY$, so that $\boxed{R \text{ lies on arc } XY\text{ (excluding the endpoints)}}$. 

 



\textbf{Solution 2}

We claim that $R$ can lie on minor arc $AB$ of the circumcircle of triangle ABC, and it can also lie on the dilation of this arc about the center of triangle $ABC$ with a factor of $-2.$

 
Let $D, E,$ and $F$ be the feet of the angle bisectors from points $A, B,$ and $C$ respectively. Trivially, $DEF$ is also the medial triangle, orthic triangle, and contact triangle (ABC is equilateral).

 
Let $I$ be the incenter of ABC. Trivially, $I$ is the centroid, orthocenter, and circumcenter of ABC (ABC is equilateral). Also, $AD=BE=CF=3r$ where $r$ is the radius of circle $\omega$ (This is trivial). $T$ is the point of tangency of $\omega$ and segment $\overline{PQ}$. 

 
R has to lie on the intersection of circles $\omega1$(center P, radius PA) and $\omega2$(center Q, radius QB), and for each choice of P, there exist two locations for R. The location that we claim to lie on the minor arc AB of the circumcircle of ABC shall be denoted M, and the other location shall be denoted N. 

 
Define triangle XYZ to be the homothety of triangle ABC about I with a factor of -2.

 
Critical claim: M, T, I, and N are collinear.

 
Proof:First we shall prove that T lies on MN using phantom points.

 
Let the intersection of MN and PQ be denoted as K. We shall prove that K and T are the same point.Let $PT = p$ and $QT = q$. Because of the equal tangent theorem, $PD=PT=p$ and $QE=QT=q$. Hence, by the pythagorean theorem (recall $AD=BE=3r$), $PA^2 = 9r^2 + p^2$ and $QB^2 = 9r^2 + q^2$. Since PN = PA and QN = QB, then $PN^2 = 9r^2 + p^2$ and $QN^2 = 9r^2 + q^2$. 

 
PQ is the perpendicular bisector of MN because MN is the radical axis of $\omega1$ and $\omega2$. Hence, M is the reflection of N across K. Also, NK is the altitude of triangle PNQ, so $PK^2-QK^2 = PN^2-QN^2 = p^2 - q^2$ by using the pythagorean theorem and earlier expressions for $PN^2$ and $QN^2$. However, $PK+QK=PQ=p+q$. Now, we have a system of equations to solve for PK and QK in terms of p and q.

 
Dividing the first equation by the second (we can do this because p+q is always nonzero), we get $PK-QK=p-q$. Combining this with our PK+QK result, we get $PK = p$ and $QK = q$. However, $PT = p$ and $QT = q$, and only one point can exist on PQ for which this result holds true. As a result, K and T are the same point, otherwise it is a contradiction. Hence M, T, and N are collinear.

 
$IT \parallel MN$. This is because both MN and IT are perpendicular to PQ (IT is perpendicular to PQ because PQ is a tangent with point of tangency T). However, both lines share point T, as discussed earlier. Hence, IT and MN are the same line, and M, T, I, and N are collinear.

 
In fact, from our earlier results from the lengths of PN and PT, we can use the pythagorean theorem to get that $NT = 3r$, a result that is always true and independent of P and Q! Also, because M is the reflection of N over K (which is the same as T), $MT = 3r$ also. However, T varies based on P and Q. On the other hand, $IT = r$ and M, T, I and N are collinear. Remembering our earlier definitions of M and N, we get that $MI = 2r$ and $IN = 4r$, with M on the opposite side of N and T from I. Hence, M can be taken to N with a homothety about I with a factor of -2, and T can be taken to M with a homothety about I with a factor of -2. Since, trivially, the circumradius of ABC is 2r (ABC is equilateral), it seems like M can lie anywhere on the circumcircle of ABC.

 
However, we must take into account the restrictions on P and Q. This limits T to only minor arc DE on the incircle of ABC, hence, because of our earlier homothety statement, M is restricted to minor arc AB on the circumcircle of ABC.  Because of our homothety statement about N, N has to lie on minor arc XY on the circumcircle of triangle XYZ.

 
Because we defined both M and N to be possible locations for R, $\fbox{R can only lie on minor arc AB of the circumcircle of triangle ABC, and also on minor arc XY of the circumcircle of triangle XYZ}$.

 
-QED

 
-Solution by thanosaops

 


	\item (\textit{USAJMO 2020}) \textbf{Note:} For any geometry problem whose statement begins with an asterisk $(*)$, the first page of the solution must be a large, in-scale, clearly labeled diagram. Failure to meet this requirement will result in an automatic 1-point deduction.


 An empty $2020 \times 2020 \times 2020$ cube is given, and a $2020 \times 2020$ grid of square unit cells is drawn  on each of its six faces. A beam is a $1 \times 1 \times 2020$ rectangular prism. Several beams are placed inside the cube subject to the following conditions:


 
                \begin{itemize}
                    \item The two $1 \times 1$ faces of each beam coincide with unit cells lying on opposite faces of the cube. (Hence, there are $3 \cdot {2020}^2$ possible positions for a beam.)
\item No two beams have intersecting interiors.
\item The interiors of each of the four $1 \times 2020$ faces of each beam touch either a face of the cube or the interior of the face of another beam.
                \end{itemize}


                What is the smallest positive number of beams that can be placed to satisfy these conditions?


 



\textbf{Solution}

Place the cube in the xyz-coordinate, with the positive x-axis pointing forward, the positive y-axis pointing right, and the positive z-axis pointing up. Let the position of a unit cube be $(x, y, z)$ if it is centered at $(x, y, z)$. Place the $2020 \times 2020 \times 2020$ cube so that the edges are parallel to the axes, and two of its corners are at $(1, 1, 1)$ and $(2020, 2020, 2020)$. Now call a beam z-oriented if its endpoints differ only in z-coordinates, and similarly call it x-oriented or y-oriented if the endpoints differ in x- or y-coordinates, respectively.

 We claim that the answer is $\boxed{3030}$. First we will prove 3030 suffices. Place the first beam so its endpoints lie at $(1, 1, 1)$, and $(2020, 1, 1)$. Place the 2nd beam so its endpoints lie at $(1, 1, 2)$ and $(1, 2020, 2)$. Place the 3rd beam so its endpoints lie at $(2, 2, 1)$ and $(2, 2, 2020)$. Now it is clear that the only faces among these three beams that are not touching the faces of the $2020 \times 2020 \times 2020$ cube of the face of another beam are the top face of beam 2, and the front and right faces of beam 3. 

 Now place the 4th beam from $(1, 3, 3)$ to $(2020, 3, 3)$, 2 units up and right of beam 1. Then place beam 5 from $(3, 1, 4)$ to $(3, 2020, 4)$, 2 units forward and up of beam 2. Place beam 6 from $(4, 4, 1)$ to $(4, 4, 2020)$, 2 unites forward and right of beam 3. Note that these 3 beams will touch the top face of beam 2, and the front and right faces of beam 3. Now the only faces that are not touching another face are the top of beam 5, and the front and right of beam 6. 

 We can repeat this process 1010 times, placing the beams in groups of 3 so that beam $3n+1$ is 2 units up and right of beam $3n-2$, beam $3n+2$ is 2 units forward and up of beam $3n-1$, and beam $3n+3$ is 2 units forward and right of beam $3n$, for all $1\leq n \leq 1009$. After placing each set of 3 as such, the faces of the previous beams will all satisfy condition 3. And for the last set of 3 (that is, beams 3028, 3029, and 3030), their faces will satisfy condition 3 as well, since the top face of beam 3029 and the front/right faces of beam 3030 will touch the cube (again, note that all other faces will be touching another beam).

 Thus, we have shown that 3030 beams suffices, so we will now show that 3030 beams is the minimum. Let there be $A$ x-oriented beams, $B$ y-oriented beams, and $C$ z-oriented beams. Now look each $2020 \times 1 \times 2020$ slice of the cube, parallel to the xz-plane. Note there are 2020 such slices. There are two cases. Either $B=2020^2$, or at least one of A or C is nonzero. In the first case, we clearly have at least $2020^2>3030$ beams as desired. In the second case, we must have at least one beam in each $2020 \times 1 \times 2020$ slice, in order to satisfy condition 3. If there is some slice without a beam, then there will be some x- or z-oriented beam with no beams touching its top or bottom face. Additionally, notice that only x- or z-oriented beams can fit in a $2020 \times 1 \times 2020$ slice, so $A+C\geq 2020$. Similarly, we must have $A=2020^2$, $C=2020^2$, or both $A+B\geq 2020$ and $B+C\geq 2020$. In every case, the number of beams, $A+B+C$, will be at least 3030 as desired.$\blacksquare$

 


	\item (\textit{USAJMO 2020}) \textbf{Note:} For any geometry problem whose statement begins with an asterisk $(*)$, the first page of the solution must be a large, in-scale, clearly labeled diagram. Failure to meet this requirement will result in an automatic 1-point deduction.


 Let $ABCD$ be a convex quadrilateral inscribed in a circle and satisfying $DA < AB = BC < CD$. Points $E$ and $F$ are chosen on sides $CD$ and $AB$ such that $BE \perp AC$ and $EF \parallel BC$. Prove that $FB = FD$.


 



\textbf{Solution}

Let $G$ be the intersection of $AE$ and $(ABCD)$ and $H$ be the intersection of $DF$ and $(ABCD)$.

 \textbf{Claim: $GH || FE || BC$}

 By Pascal's on $GDCBAH$, we see that the intersection of $GH$ and $BC$, $E$, and $F$ are collinear. Since $FE || BC$, we know that $HG || BC$ as well. $\blacksquare$

 Note that since all cyclic trapezoids are isosceles, $HB = GC$. Since $AB = BC$ and $EB \perp AC$, we know that $EA = EC$, from which we have that $DGCA$ is an isosceles trapezoid and $DA = GC$. It follows that $DA = GC = HB$, so $BHAD$ is an isosceles trapezoid, from which $FB = FD$, as desired. $\blacksquare$

 



\textbf{Solution 2}

Let $G=\overline{FE}\cap\overline{BC}$, and let $G=\overline{AC}\cap\overline{BE}$. Now let $x=\angle ACE$ and $y=\angle BCA$.

 From $BA=BC$ and $\overline{BE}\perp \overline{AC}$, we have $AE=EC$ so $\angle EAC =\angle ECA = x$. From cyclic quadrilateral ABCD, $\angle ABD = \angle ACD = x$. Since $BA=BC$, $\angle BCA = \angle BAC = y$.

 Now from cyclic quadrilateral ABC and $\overline{FE}\parallel \overline{BC}$ we have $\angle FAC = \angle BAC = \pi - \angle BCD = \pi - \angle FED$. Thus F, A, D, and E are concyclic, and $\angle DFG = \angle DAE = \angle DAC - \angle EAC = \angle DBC - x$ Let this be statement 1.

 Now since $\overline{AH}\perp \overline {BH}$, triangle ABC gives us $\angle BAH + \angle ABG = \frac{\pi}{2}$. Thus $y+x+\angle GBE=\frac{\pi}{2}$, or $\angle GBE = \frac{\pi}{2}-x-y$. 

 Right triangle BHC gives $\angle HBC = \frac{\pi}{2}-y$, and $\overline{BC}\parallel \overline{FE}$ implies $\angle BEG=\angle HBC = \frac{\pi}{2}-y.$ 

 Now triangle BGE gives $\angle BGE = \pi - \angle BEG - \angle GBE = \pi - (\frac{\pi}{2}-y)-(\frac{\pi}{2}-x-y)=x+2y$. But $\angle FGB = \angle BGE$, so $\angle FGB=x+2y$. Using triangle FGD and statement 1 gives 
 \begin{align*}\angle FDG &= \pi - \angle DFG - \angle FGB \\ &= \pi - (\angle DBC - x) - (x + 2y) \\ &= \pi - (\angle GBE + \angle EBC - x) - (x + 2y) \\ &= \pi - ([\frac{\pi}{2}-x-y]+[\frac{\pi}{2}-y]-x)-(x+2y) \\ &= x \\ &= \angle FBD\end{align*}


 Thus, $\angle FDB = \angle FBD$, so $\boxed{FB=FD}$ as desired.$\blacksquare$

 



\textbf{Solution 3 (Angle-Chasing)}

Proving that $FB=FD$ is equivalent to proving that $\angle FBD= \angle FDB$. Note that $\angle FBD=\angle ACD$ because quadrilateral $ABCD$ is cyclic. Also note that $\angle BAC=\angle ACB$ because $AB=BC$. $AE=EC$, which follows from the facts that $BE \perp AC$ and $AB=AC$, implies that $\angle CAE= \angle ACE= \angle ACD= \angle FBD$. Thus, we would like to prove that triangle $FBD$ is similar to triangle $AEC$. In order for this to be true, then $\angle BFD$ must equal $\angle AEC$ which implies that $\angle AFD$ must equal $\angle AED$. In order for this to be true, then quadrilateral $AFED$ must be cyclic. Using the fact that $EF \parallel BC$, we get that $\angle AFE= \angle ABC$, and that $\angle FED= \angle BCE$, and thus we have proved that quadrilateral $AFED$ is cyclic. Therefore, triangle $FDB$ is similar to isosceles triangle $AEC$ from AA and thus $FB=FD$.

 -xXINs1c1veXx

 



\textbf{Solution 4}

BE is perpendicular bisector of AC, so $\angle ACE = \angle EAC$. FE is parallel to BC and ABCD is cyclic, so AFED is also cyclic. $\angle AFD = \angle AED = 2 \angle ACE = 2 \angle FBD$. Hence, $\angle FBD = \angle BDF$, $FD = FB$. 

 Mathdummy

 



\textbf{Solution 5}

Let $G$ be on $AB$ such that $GD \perp DB$, and $H = GD \cap EB$. Then $\angle{ADB} = \angle{EDB} = 90^{\circ} - \angle{ABE} \implies \triangle{DAE}$ is the orthic triangle of $\triangle{HGB}$. Thus, $F$ is the midpoint of $GB$ and lies on the $\perp$ bisector of $DB$.

 



\textbf{Solution 6}

Let $FE$ meet $AC$ at $J$, $BE$ meet $AC$ at $S$, connect $AE, SD$.Denote that $\angle{BCA}=\alpha; AB=BC, \angle{BAC}=\angle{BCA}=\alpha$, since $EF$ is parallel to $BC$, $\angle{AJF}=\angle{ACB}=\alpha$. $\angle{AJF}$and $\angle{EJS}$ are vertical angle, so they are equal to each other. $BE\bot{AC}$,$\angle{JES}=90^{\circ}-\alpha$, since $\angle{EFB}=\angle{AJF}+\angle{FAJ}=2\alpha$, we can express $\angle{FBE}=180^{\circ}-2\alpha-(90^{\circ}-\alpha)=90^{\circ}-\alpha= \angle{FEB}$, leads to $FE=FB$

 Notice that quadrilateral $AFED$ is a cyclic quadrilateral since $\angle{ADE}+\angle{AFE}=\angle{ADE}+\angle{ABC}=180^{\circ}$.

 Assume $\angle{ECA}=\beta$, $\triangle{AES}$ is congruent to $\triangle{CES}$ since $AS=AS,\angle{ASE}=\angle{BSE}, SE=SE(SAS)$, so we can get $\angle{EAS}=\beta$ Let the circumcircle of $AFED$ meets $AC$ at $Q$Now notice that $\widehat{QE}=\widehat{QE}, \angle{QAE}=\angle{QDE}=\beta$; similarly, $\widehat{FQ}=\widehat{FQ}; \angle{FDQ}=\angle{FAQ}=\alpha$.

 $\angle{FDE}=\alpha+\beta; \angle{FED}=\angle{BCD}=\alpha+\beta$, it leads to $FD=FE$. 

 since $FE=FB;FD=FE, DF=BF$ as desired~bluesoul

 


	\item (\textit{USAJMO 2020}) \textbf{Note:} For any geometry problem whose statement begins with an asterisk $(*)$, the first page of the solution must be a large, in-scale, clearly labeled diagram. Failure to meet this requirement will result in an automatic 1-point deduction.


 Suppose that $(a_1,b_1),$ $(a_2,b_2),$ $\dots,$ $(a_{100},b_{100})$ are distinct ordered pairs of nonnegative integers. Let $N$ denote the number of pairs of integers $(i,j)$ satisfying $1\leq i<j\leq 100$ and $|a_ib_j-a_jb_i|=1$. Determine the largest possible value of $N$ over all possible choices of the $100$ ordered pairs.


 



\textbf{Solution}

Call the pair $(i, j)$ good if $1\leq i < j \leq 100$ and $|a_ib_j-a_jb_i|=1$. Note that we can reorder the pairs $(a_1, b_1), (a_2, b_2), \ldots, (a_{100}, b_{100})$ without changing the number of good pairs. Thus, we can reorder them so that $a_1\leq a_2\leq\ldots\leq a_{100}$. Furthermore, reorder them so that if $a_i=a_j$ for some $i<j$, then $b_i<b_j$. 

 Now I claim the maximum value of $N$ is $\boxed{197}$. First, we will show $N\leq 197$.

 The key claim is that for any fixed integer $m$, $1\leq m \leq 100$, then there are at most 2 values of $i$ with $1\leq i<m$ with the pair $(i, m)$ good. We will prove this with contradiction.

 Assume there are 3 integers $i$, $j$, and $k$ that are less than $m$, such that $(i, m)$, $(j, m)$, and $(k, m)$ are all good. Note that $gcd(a_m, b_m)$ divides $a_ib_m-a_mb_i=\pm 1$, so $gcd(a_m, b_m)=1$. From $a_ib_m-a_mb_i=\pm 1$, we have 
 \[a_ib_m\equiv \pm 1 \mod a_m\]Similarly, 
 \[a_jb_m\equiv \pm 1 \mod a_m\] 
 \[a_kb_m\equiv \pm 1 \mod a_m\]

 From $gcd(a_m, b_m)=1$ we have 
 \[a_i\equiv \pm a_j \equiv \pm a_k \mod a_m\]

 But since $i$, $j$, and $k$ are less than $m$, from our reordering we have $0\leq a_i<a_m$, $0\leq a_j<a_m$, and $0\leq a_k<a_m$. Thus our congruence gives us that $a_j=a_i$ or $a_j=a_m-a_i$, and similar for $a_k$. Thus at least two of $a_i, a_j, a_k$ are equal. WLOG let them be $a_i$ and $a_j$. Now we have 4 cases.

 \textbf{Case 1: $a_m>2$}

 We have $a_ib_m-a_mb_i=\pm 1$ and $a_jb_m-a_mb_j=\pm 1$, so since $a_i=a_j$ we can subtract the two equations to get $a_mb_i-a_mb_j$ equals $\pm 2, 0$. But $a_m>2$, so $|a_m(b_i-b_j)|\leq 2$ implies that $b_i=b_j$. But all pairs are distinct and $a_i=a_j$, so this is a contradiction. 

 \textbf{Case 2: $a_m=2$}

 We have 2 subcases.

 \textbf{Subcase 2a: $a_i\equiv a_j\equiv a_k\equiv 1 \mod a_m$}

 Clearly we must have $a_i=a_j=a_k=1$ in this case, as all 3 variables are less than $a_m=2$. Then $1\cdot b_m-a_mb_i=\pm 1$, $1\cdot b_m-a_mb_j=\pm 1$, and $1\cdot b_m-a_mb_k=\pm 1$. So at least 2 of $b_i$, $b_j$, and $b_k$ must be equal, but this contradicts the fact that all pairs $(a_x, b_x)$ are distinct.

 \textbf{Subcase 2b: $a_i\equiv a_j\equiv a_k\equiv 0 \mod a_m$}

 $a_i$ must be even in this case, but since $a_m$ is also even then $a_ib_m-a_mb_i$ must be even, which contradicts $|a_ib_m-a_mb_i|=1$.

 \textbf{Case 3: $a_m=1$}

 We must have $a_i=0,1$, and similarly for $a_j$ and $a_k$. Thus at least 2 of the 3 must be equal. Now we have 2 subcases again.

 \textbf{Subcase 3a: 2 of the 3 are equal to 0}

 WLOG let $a_i=a_j=0$. Then $a_ib_m-a_mb_i=\pm 1$ and $a_m=1$ so $b_i=\pm 1$. Since $b_i$ is nonnegative, $b_i=1$. Similarly $b_j=1$, but all pairs are distinct so we have a contradiction.

 \textbf{Subcase 3b: 2 of the 3 are equal to 1}

 WLOG let $a_i=a_j=0$. Then $a_ib_m-a_mb_i=\pm 1$, so $b_m-b_i=\pm 1$. But $i<m$ and $a_i=a_m=1$, so by our ordering of the pairs we have $b_i<b_m$ so $b_m-b_i=1$. But similarly $b_m-b_j=1$, so $b_i=b_j$, which is a contradiction to all pairs being distinct.

 \textbf{Case 4: $a_m=0$}

 We must have $a_i=0$, since $0\leq a_i\leq a_m=0$. Thus $a_ib_m-a_mb_i=0$, but $|a_ib_m-a_mb_i|=1$, contradiction.$\blacksquare$

 Thus our claim is proved. Now for $m=1$, there is clearly no $i$ with $1\leq i < m \leq 100$ and $(i, m)$ good. For $m=2$, there is at most one value of i with $1\leq i < m \leq 100$ and $(i, m)$ good (namely, this is when $i=1$). By our claim, for all $2<m\leq 100$ there are at most 2 values of $i$ with $1\leq i < m \leq 100$ and $(i, m)$ good, so adding these up gives at most $2\cdot 98+1=197$ good pairs. Now we will show that 197 good pairs is achievable. 

 The sequence $(1,1), (1,2), (2,3),\ldots ,(99,100)$ has 197 good pairs - namely, $(1, n)$ is good for all $1<n\leq 100$, and $(i, i+1)$ is good for all $1<i<100$. This adds up for a total of 197 good pairs, so we are done. $\blacksquare$

 



\textbf{Solution 2}

We claim the answer is $2n-3$ (so $197$), where $100$ is to be replaced with $n$.

 The problem essentially counts unordered pairs $(P,Q)$ of lattice points so that $POQ$ has area $\frac{1}{2}$ due to Shoelace ($O$ is the origin).

 For a construction, take the points $(0,1),(1,1),(1,2), \dots (1,n-1)$.

 We now show by induction that $2n-3$ is the maximum, where $n=2$ obviously holds as a base case.

 To reduce to the inductive hypothesis of $n-1$ points from $n$ points, we are allowed to consider the point with the greatest sum of coordinates $R$.

 The key claim is that $R$ is part of at most $2$ pairs (triangles) with area $\frac{1}{2}$.

 Before we proceed, note that $R$ must have relatively prime coordinates $(m,n)$, or else any triangles found can be reduced down by a greatest common factor along $OR$, impossible since $\frac{1}{2}$ is the minimum positive area.

 Call the hypothetical point(s) $X$. With triangle $XOR$ with area $\frac{1}{2}$, we know the length of $OR$, so indeed the locus of $X$ is two lines evenly spaced next to $OR$ each parallel to $OR$ (before we consider the lattice condition).

 When we add in the lattice condition, we see that each of these two lines contain at most one valid $X$. This is because if two lattice points with non-negative integers lied on one of these lines, one would have to be a multiple of $m$ more than the other on the $x$-coordinates and a multiple of $n$ more on the $y$-coordinates. This contradicts the coordinate sum maximality of $R$.

 Thus, by deleting $R$, we lose at most two pairs $(X,R)$, completing the induction and the problem.

 


	\item (\textit{USAJMO 2020}) \textbf{Note:} For any geometry problem whose statement begins with an asterisk $(*)$, the first page of the solution must be a large, in-scale, clearly labeled diagram. Failure to meet this requirement will result in an automatic 1-point deduction.


 Let $n \geq 2$ be an integer. Let $P(x_1, x_2, \ldots, x_n)$ be a nonconstant $n$-variable polynomial with real coefficients. Assume that whenever $r_1, r_2, \ldots , r_n$ are real numbers, at least two of which are equal, we have $P(r_1, r_2, \ldots , r_n) = 0$. Prove that $P(x_1, x_2, \ldots, x_n)$ cannot be written as the sum of fewer than $n!$ monomials. (A monomial is a polynomial of the form $cx^{d_1}_1 x^{d_2}_2\ldots x^{d_n}_n$, where $c$ is a nonzero real number and $d_1$, $d_2$, $\ldots$, $d_n$ are nonnegative integers.)


 


\end{enumerate}